\phantomsection\label{fig:vol04:mongol-conquest:empire-formation}
\begin{center}
\begin{tikzpicture}[x=.88cm,y=.88cm,font=\small]
  \fill[paper] (0,0) rectangle (11.8,6.45);

  \fill[source!16] (.55,3.75) -- (5.35,3.55) -- (5.2,5.95) -- (.75,5.85) -- cycle;
  \node[align=center] at (2.85,4.95) {草原政治中心\\诸王、军队与继承};
  \fill[timeline!20] (5.75,3.65) -- (11.2,3.8) -- (11.15,5.85) -- (5.9,5.8) -- cycle;
  \node[align=center] at (8.65,4.95) {西北与中亚\\多方向扩张的战区};
  \fill[dynasty!19] (.55,.65) -- (11.15,.75) -- (11.15,3.25) -- (.6,3.4) -- cycle;
  \node[align=center] at (8.95,1.45) {金、南宋、大理与江南\\城池、道路、水系和接收网络};

  \node[draw=source,fill=paper,rounded corners=2pt,align=center,inner sep=3pt] (qaraqorum)
    at (2.8,4.38) {1206\\蒙古诸部统一};
  \node[draw=timeline,fill=paper,rounded corners=2pt,align=center,inner sep=3pt] (xia)
    at (4.8,3.0) {1209、1227\\西夏};
  \node[draw=dynasty,fill=paper,rounded corners=2pt,align=center,inner sep=3pt] (jin)
    at (6.35,2.45) {1211--1234\\金与蔡州};
  \node[draw=timeline,fill=paper,rounded corners=2pt,align=center,inner sep=3pt] (centralasia)
    at (8.55,4.35) {1218--1219\\西辽、花剌子模};
  \node[draw=dynasty,fill=paper,rounded corners=2pt,align=center,inner sep=3pt] (dali)
    at (4.45,1.12) {1253\\大理};
  \node[draw=dynasty,fill=paper,rounded corners=2pt,align=center,inner sep=3pt] (xiangyang)
    at (8.35,2.2) {1268--1273\\襄樊};
  \node[draw=dynasty,fill=paper,rounded corners=2pt,align=center,inner sep=3pt] (jiangnan)
    at (10.0,1.1) {1276--1279\\临安、崖山};

  \draw[->,ultra thick,color=source!85] (qaraqorum) -- node[left,font=\scriptsize]{西夏战事} (xia);
  \draw[->,ultra thick,color=timeline!85] (qaraqorum) -- node[above,sloped,font=\scriptsize]
    {西北远征} (centralasia);
  \draw[->,ultra thick,color=dynasty!85] (xia) -- node[above,sloped,align=center,font=\scriptsize]
    {华北战争与\\宋蒙短期协作} (jin);
  \draw[->,thick,color=dynasty!85] (jin) -- node[below,sloped,font=\scriptsize]
    {西南战略环境改变} (dali);
  \draw[->,ultra thick,color=dynasty!85] (jin) -- node[above,sloped,align=center,font=\scriptsize]
    {围城、水陆补给\\与降附者} (xiangyang);
  \draw[->,ultra thick,color=dynasty!85] (xiangyang) -- node[below,sloped,font=\scriptsize]
    {南下与江海战场} (jiangnan);

  \node[draw=source!75,fill=paper,rounded corners=2pt,align=center,inner sep=3pt] at (2.4,2.05)
    {征服后的难题：\\户籍、赋税、工匠、道路\\与地方军政的接收};
  \node[text=source,align=center,font=\scriptsize] at (6.1,.75)
    {箭头表示相互关联的战区与阶段，不表示单一路线、行军距离或同时控制的疆域。};
\end{tikzpicture}

\smallskip
\small\textit{图\quad 蒙古扩张至元朝形成的若干战区与接收节点（1206--1279，示意）：据《元朝秘史》《元史》〈太祖本纪〉、〈宪宗本纪〉、〈世祖本纪〉及《新元史》相关编年材料，并与《金史》〈哀宗本纪〉、《宋史》〈理宗本纪〉、〈度宗本纪〉和《续资治通鉴》互读。图中列出西夏、金、中亚、大理、襄樊和江南等相连而不等同的战区，提示军事推进须同继承政治、围城补给及征服后的行政接收一起理解；路线、范围、城池控制和各方兵力均为约略示意。}
\end{center}
